\documentclass[10pt,a4paper]{amsart}
\usepackage[margin=19mm]{geometry}
\usepackage{amsmath,amssymb,mathtools}
\usepackage[scr=rsfs,cal=euler]{mathalfa}
\usepackage{xfrac}
\usepackage[unicode,hidelinks,bookmarks=false]{hyperref}
\newcommand{\ie}{i.e.\ }
\newcommand{\cf}{cf.\ }
\newcommand{\resp}{respectively\ }
\newcommand{\Subsection}{\subsection}
\providecommand{\nobreakdash}{\nobreak\hbox{-}}
\setlength{\emergencystretch}{2em}
\multlinegap0pt
\usepackage{amssymb}
\usepackage{xcolor}
\usepackage{mathtools}
\usepackage[shortlabels]{enumitem}
\usepackage{bm}
\newtheorem{lemma}{Lemma}
\newcommand{\E}{\mathbb{E}}
\def\bbE{\mathbb E}
\def\tr{{\rm tr}}
\title{Intermittency of geometric Brownian motion on \textbf{SL}(\lowercase{n})}
\author{Sefika Kuzgun, Felix Otto, Christian Wagner}
\date{Source: arXiv:2602.07167v2 (CC BY 4.0)}
\begin{document}
\maketitle

\noindent\emph{Source and licence.} Intermittency of geometric Brownian motion on \textbf{SL}(\lowercase{n}). Version arXiv:2602.07167v2 (CC BY 4.0), distributed by arXiv under the Creative Commons Attribution 4.0 International licence, http://creativecommons.org/licenses/by/4.0/, as recorded in the arXiv licence metadata for this exact version and shown on the abstract page https://arxiv.org/abs/2602.07167v2. Full licence notice: \texttt{licenses/arxiv-2602.07167-CC-BY-4.0.txt}.\par\smallskip

\noindent\textit{Editorial scope.} Counted unit: the article's lemma lem:zeta (source0.tex lines 688--696). The article's complete original proof of it is delivered with it (source0.tex lines 927--1004, one \texttt{proof} environment of 78 lines, which the article places away from the statement). No mathematics is written, repaired or re-derived: the statement and the proof are the article's own lines, in the article's own notation, and no retained line is substituted or re-flowed.  Retained uncounted as the source-local context the proof needs: lines 485--488; lines 517--519; lines 540--542; lines 583--621; lines 624--626; lines 627--629; lines 635--636; lines 643--645; lines 647--649; lines 656--658; lines 669--672; lines 682--684.  Nothing was omitted from any retained span, so the package carries no omission record.  No retained line contains a citation, so the wrapper carries no bibliography and no bibliography entry.  No retained line contains a figure environment, an image-inclusion command or drawing code, so no image is delivered and the package declares no asset dependency.  Editorial formatting only, in addition: an AMS \texttt{amsart} preamble that restates the article's own declarations and macros used by the retained lines, namely the environments \texttt{lemma} on separate counters, together with the standard packages those lines need.  The retained environments print in this document as Lemma 1 in order of appearance, whereas the article prints the same items with the numbering of its own full document.  The complete native source of the article as distributed in the arXiv e-print for this exact version is bundled unchanged as \texttt{sources/194154/source0.tex}.
 \begin{align}
    d\tr G &=\tr G_\tau (\circ dB+\circ dB^*) =2\tr G_\tau \circ dB_{sym}, \label{trgstr} \\
    d\tr G^2&=2\tr G^2_\tau (\circ dB+\circ dB^*) =4\tr G^2_\tau\circ dB_{sym}. \label{trG2str} 
\end{align}

Hence the Itô formulation of the equation \eqref{trgstr} becomes \begin{align}
    d\tr G &=2\tr G_\tau dB_{sym}+\tr G_\tau d\tau. \label{trg} 
\end{align} 

\begin{align}\label{qv}
    d[\tr  G B_{ sym} \tr  G B_{ sym}]=\frac{1}{\alpha_n} (n \tr  G_\tau^2-\tr^2 G_\tau )d\tau.
\end{align} 

\begin{lemma}\label{Lem:2by2}
  It holds
   %
   \begin{align*}
    \frac{d}{d\tau}\begin{bmatrix}
        \bbE \tr^2 G \\ \bbE \tr G^2
    \end{bmatrix}  = A \begin{bmatrix}
        \bbE \tr^2 G_\tau \\ \bbE \tr G_\tau^2
    \end{bmatrix}, \quad \quad
   \begin{bmatrix}
        \bbE \tr^2 G_{\tau=0} \\ \bbE \tr  G_{\tau=0}^2
    \end{bmatrix} = \begin{bmatrix}
    n^2\\n
\end{bmatrix}
\end{align*}
%
where 
%
\begin{align*}
    A=(2-\frac{4}{\alpha_n}) \begin{bmatrix}
        1&0\\
        0&1
    \end{bmatrix}+\frac{2n}{\alpha_n}\begin{bmatrix}
        0&2\\
        1&1
    \end{bmatrix}.
\end{align*}
%
The matrix $A$ has the eigenvalues $\lambda_1=2+\frac{4}{n+2}> \lambda_2=2-\frac{2}{n-1}$ with left eigenvectors $v_1=\begin{bmatrix}
    1\\2
\end{bmatrix}$ and $v_2=\begin{bmatrix}
    1\\-1
\end{bmatrix}$. As a consequence, for all $\tau>0$, 
%
\begin{align}\label{difference}
 \bbE  \tr^2G_{\tau}-\bbE \tr G_\tau^2  = n(n-1) e^{\lambda_2\tau}.  
\end{align}
%
\end{lemma}

\begin{align}\label{R}
    R_\tau:=\tr G_\tau.
\end{align}

Indeed, recall from \eqref{trg} and \eqref{qv} that $R$ satisfies  \begin{align}
dR=2\tr G_\tau dB_{sym}+R_\tau d\tau, \quad \quad R_{\tau=0}=n ,\label{RSDE}
\end{align}

\begin{align}\label{QuadVarR}
        d[RR]&=\frac{4}{n+2}R_\tau^2 d\tau+\frac{4n}{\alpha_n} (\tr G_\tau^2-\tr^2 G_\tau)d\tau .   \end{align}

\begin{align}\label{Rhat}
  R=n e^{\tau} \hat R ,
\end{align}

 \begin{align}\label{Rhatapprox}
   d\hat R \approx \frac{2}{\sqrt{n+2}}\hat R_{ \tau} dw, \quad \quad \hat R_{ \tau=0}=1. 
\end{align}

\begin{align}\label{backwardpde}
	   \frac{\partial \zeta}{\partial \tau}+r \frac{\partial \zeta}{\partial r}+\frac{2 }{n+2}r^2\frac{\partial^2\zeta}{\partial  r^2}=0.
\end{align}

 \begin{align}
    \zeta(\tau,\hat r)=\hat r \hat 
\zeta(\tau,\hat r) \label{zetatohatzeta}
\end{align} 

\begin{align}\label{sigmahat}
    \hat \sigma = \ln \hat r ,
\end{align}

\begin{lemma}\label{lem:zeta} 
   Suppose $\hat \zeta=\hat\zeta(\tau,\hat\sigma)$ solves \begin{align}\label{zetahat2}
\frac{\partial \hat \zeta}{\partial \tau}+\frac{2}{n+2}\frac{\partial \hat \zeta}{\partial\hat \sigma }+ \frac{2}{n+2}\frac{\partial^2 \hat \zeta}{\partial \hat \sigma^2}=0,
\end{align}
then \begin{align} \label{zetahatineq}
    \bbE \hat R \hat \zeta(\tau, \hat R) - \hat \zeta(\tau=0,\hat r= 1 ) \lesssim_{ n } \int_0^{\tau} d\tau' e^{-(1-\frac{\lambda_2}{2})\tau'} \sup_{\hat \sigma} \left|\frac{\partial \hat \zeta}{\partial\hat \sigma}+ \frac{\partial^2 \hat \zeta}{\partial\hat \sigma^2} \right|  ,
 \end{align}
 where $\lambda_2$ is defined in Lemma~\ref{Lem:2by2}. 
\end{lemma}

\begin{proof}[Proof of Lemma~\ref{lem:zeta}] We start from \eqref{RSDE} and \eqref{QuadVarR} that imply for $ \zeta = \zeta ( \tau, R_{ \tau } ) $
%
\begin{align}\label{eqn02} \nonumber
	d \zeta 
	%
	&= \text{martingale} + \Big( \frac{ \partial \zeta }{ \partial \tau } +R_{ \tau }  \frac{ \partial \zeta }{ \partial r } \Big) d \tau 
	+ \frac{ 1 }{ 2 } \frac{ \partial^2 \zeta }{ \partial r^2 } d [ R \, R ] \\
	%
	&= \text{martingale} + \Big( \frac{ \partial \zeta }{ \partial \tau } + R_{ \tau } \frac{ \partial \zeta }{ \partial r } + \frac{ 2 }{ n + 2 } R_{ \tau }^2 \frac{ \partial^2 \zeta }{ \partial r^2 }  \Big) d \tau
	+ \frac{ 2 n }{ \alpha_{ n } } \Big( \tr G_{ \tau }^2 - \tr^2 G_{ \tau } \Big) \frac{ \partial^2 \zeta }{ \partial r^2 } d \tau .
\end{align}
%
Due to the term $\tr  G^2$, this differential equation \eqref{eqn02} even on the level of expectation is not closed. To overcome this, we will invoke Lemma~\ref{Lem:2by2}, more specifically identity \eqref{difference}. 
%

\noindent
Choosing our observable $ \zeta $ in \eqref{eqn02} according to \eqref{backwardpde} as we argued to be equivalent to \eqref{zetahat2} the sum in the first parenthesis vanishes. Indeed, recall from \eqref{Rhat}, \eqref{zetatohatzeta} and \eqref{sigmahat} that we change variables according to
%
\begin{align*}
	r = n e^{ \tau } \hat r ,
	\quad \zeta = \hat r \, \hat \zeta ,
	\quad \hat \sigma = \ln \hat r ,
\end{align*}
%
so that $ \frac{ \partial }{ \partial r } = e^{ - \tau } \frac{ 1 }{ n } \frac{ \partial }{ \partial \hat r }  $, $ \frac{ \partial }{ \partial \hat r } = \frac{ 1 }{ \hat r } \frac{ \partial }{ \partial \hat\sigma } $ and readily
%
\begin{align}\label{Trafo}
	\frac{ \partial \zeta }{ \partial \tau } = \hat r \frac{ \partial \hat \zeta }{ \partial \tau } - \hat r \Big( \hat \zeta + \frac{ \partial \hat \zeta }{ \partial \hat \sigma } \Big)  ,
	\quad r \frac{ \partial \zeta }{ \partial r } = \hat r \Big( 1 + \frac{ \partial  \hat \zeta }{ \partial \hat\sigma } \Big) , 
	\quad  r^2 \frac{ \partial^2 \zeta }{ \partial r ^2} = \hat r \Big( \frac{ \partial \hat \zeta }{ \partial \hat\sigma } + \frac{ \partial^2 \hat \zeta }{ \partial \hat\sigma^2 } \Big) .
\end{align}
%
Thus, for the term of our interest
%
\begin{align*}
	\frac{ \partial \zeta }{ \partial \tau } + r \frac{ \partial \zeta }{ \partial r }  + \frac{ 2 }{ n + 2 } r^2\frac{ \partial^2 \zeta }{ \partial r^2 }
	\stackrel{(\ref{Trafo})}{ = } \hat r \Big( \frac{ \partial \hat \zeta }{ \partial \tau } + \frac{ 2 }{ n + 2 } \Big( \frac{ \partial \hat \zeta }{ \partial \hat\sigma } + \frac{ \partial^2 \hat\zeta }{ \partial \hat\sigma^2 } \Big) \Big) ,
\end{align*}
%
which makes the operator in \eqref{zetahat2} appear, and \eqref{eqn02} implies the preliminary
%
\begin{align}\label{eqn01}
	\frac{d  \E \zeta }{ d \tau }
	%
	\lesssim_n \sup_{ \hat \sigma } \Big| r \frac{ \partial^2 \zeta }{ \partial r^2 } \, \Big|
	\, \E \frac{ 1 }{ R_\tau } \big( \tr^2 G_{ \tau } - \tr G_{ \tau }^2 \big)  
	%
	\stackrel{ (\ref{Trafo}) }{ \approx_n } \sup_{ \hat \sigma } e^{ - \tau } \Big| \frac{ \partial \hat \zeta }{ \partial \hat\sigma } + \frac{ \partial^2 \hat \zeta }{ \partial \hat\sigma^2 }
	 \Big|
	\, \E \frac{ 1 }{ R_{ \tau } } \big( \tr^2 G_{ \tau } - \tr G_{ \tau }^2 \big)  .
\end{align}
%
On the remaining expectation we use the elementary inequality
%
\begin{align*}
	0 \leq \frac{ 1 }{ R_{ \tau } } \big( \tr^2 G_{ \tau } - \tr G_{ \tau }^2 \big)
	\stackrel{ (\ref{R}) }{ = } \frac{ 1 }{   \tr G_{ \tau } } \big( \tr^2 G_{ \tau } - \tr G_{ \tau }^2 \big) 
	\leq \sqrt{ \tr^2 G_{ \tau } - \tr G_{ \tau }^2 } ,
\end{align*}
%
so that by Jensen's inequality and Lemma \ref{Lem:2by2}
%
\begin{align*}
	\E \frac{ 1 }{ R_{ \tau } } \big( \tr^2 G_{ \tau } - \tr G^2_{ \tau } \big)
	\leq \sqrt{ \E ( \tr^2 G_{ \tau } - \tr G_{ \tau }^2 ) } 
	\stackrel{ (\ref{difference}) }{ \lesssim_n } e^{ \frac{ \lambda_2 }{ 2 } \tau } .
\end{align*}
%
Hence \eqref{eqn01} turns into:
%
\begin{align*}
	\frac{d  \E \hat R \hat \zeta }{ d \tau } 
	%
	\lesssim_n e^{ - ( 1 - \frac{ \lambda_2 }{ 2 } ) \tau } \sup_{ \hat \sigma } \Big| \frac{ \partial \hat \zeta }{ \partial \hat\sigma } + \frac{ \partial^2 \hat \zeta }{ \partial \hat\sigma^2 }
	\Big| .
\end{align*}
%
Integrating the estimate and using the second item in \eqref{Rhatapprox} yields the claim in \eqref{zetahatineq}.\hfill\end{proof}

\end{document}
